%------------------------
% Resume Template
% Author : Anubhav Singh
% Github : https://github.com/xprilion
% License : MIT
%------------------------

\documentclass[a4paper,20pt]{article}

\usepackage{latexsym}
\usepackage[empty]{fullpage}
\usepackage{titlesec}
\usepackage{marvosym}
\usepackage[usenames,dvipsnames]{color}
\usepackage{verbatim}
\usepackage{enumitem}
\usepackage{tabularx}
\usepackage[pdftex]{hyperref}
\usepackage{fancyhdr}
\usepackage[utf8]{inputenc}
\usepackage[T1]{fontenc}
\usepackage{ragged2e}

\pagestyle{fancy}
\fancyhf{} % clear all header and footer fields
\fancyfoot{}
\renewcommand{\headrulewidth}{0pt}
\renewcommand{\footrulewidth}{0pt}

% Adjust margins
\addtolength{\oddsidemargin}{-0.530in}
\addtolength{\evensidemargin}{-0.375in}
\addtolength{\textwidth}{1in}
\addtolength{\topmargin}{-.45in}
\addtolength{\textheight}{1in}

\urlstyle{rm}

\raggedbottom
\raggedright
\setlength{\tabcolsep}{0in}

% Sections formatting
\titleformat{\section}{
  \vspace{-4pt}\bfseries\raggedright\large
}{}{0em}{}[\color{black}\titlerule \vspace{-3pt}]

%-------------------------
% Custom commands
\newcommand{\resumeItem}[2]{
  \item\small\justifying{
    \textbf{#1}{#2}
  }
}

\newcommand{\resumeSubheading}[4]{
  \item
    \begin{tabular*}{\linewidth}{l@{\extracolsep{\fill}}r}
      \textbf{#1} & #2 \\
      \textsc{#3} & \textit{#4} \\
    \end{tabular*}
}


\newcommand{\resumeSubItem}[2]{\resumeItem{#1}{#2}\vspace{-3pt}}

% Linha de competências por experiência (sem marcador, pequena, rótulo em negrito)
\newcommand{\resumeSkills}[1]{\item[]\small{\textbf{Competências:}~#1}}

% Publicação: coluna de ano + título em negrito + veículo (estilo da referência)
\newcommand{\resumePub}[3]{%
  \noindent\begin{tabularx}{\textwidth}{@{}p{2.4em}@{\hspace{0.4em}}X@{}}
    {\small\textcolor[gray]{0.4}{#1}} & {\small\textbf{#2}, #3}
  \end{tabularx}\vspace{1pt}
}

% Conferência: coluna de ano + conferência, duração + local (à direita)
\newcommand{\resumeConf}[4]{%
  \noindent\begin{tabularx}{\textwidth}{@{}p{2.4em}@{\hspace{0.4em}}X@{}r@{}}
    {\small\textcolor[gray]{0.4}{#1}} & {\small #2, #3} & {\small\textit{#4}}
  \end{tabularx}\vspace{1pt}
}

\renewcommand{\labelitemii}{$\circ$}

\newcommand{\resumeSubHeadingListStart}{\begin{itemize}[leftmargin=*, topsep=0pt, itemsep=6pt, parsep=0pt]}
\newcommand{\resumeSubHeadingListEnd}{\end{itemize}}
\newcommand{\resumeItemListStart}{\begin{itemize}[leftmargin=*, topsep=1pt, itemsep=1pt, parsep=0pt]}
\newcommand{\resumeItemListEnd}{\end{itemize}}

%-----------------------------
%%%%%%  CV STARTS HERE  %%%%%%


\begin{document}

%----------HEADING-----------------
\textbf{{\LARGE Arthur Cancellieri Pires}} \\ \vspace{2pt}
\textit{Líder em Ciência de Dados \textbar{} Python, PyTorch \textbar{} Detecção de Anomalias} \\ \vspace{3pt}
{\small \href{mailto:arthur.cancellieri.pires@gmail.com}{arthur.cancellieri.pires@gmail.com}
 \textbar{} \href{https://www.linkedin.com/in/arthur-cancellieri-pires-730963160}{LinkedIn: Arthur Cancellieri Pires}
 \textbar{} \href{https://github.com/pirao}{GitHub: @Pirao}
 \textbar{} Vitória, Brasil}

\vspace{-8pt}

%----------Resumo Profissional-----------------
\section{Resumo Profissional}
\vspace{1pt}
{\small\justifying Líder em Ciência de Dados com mais de 7 anos transformando dados de sensores industriais em ML de produção para manutenção preditiva. Especialista em Python, PyTorch e deep learning para séries temporais e detecção de anomalias, entregando previsões de falha em tempo real com 91\% de precisão que reduziram paradas não planejadas e custos de manutenção. Liderou uma equipe de 12 pessoas e levou ML à produção de ponta a ponta: pipelines escaláveis, deploy de modelos e APIs (FastAPI, Docker). Doutorando com mestrado em engenharia mecânica, unindo rigor de pesquisa à produção real.\par}

\vspace{-6pt}

%----------Work Experience-----------------

\section{Experiência Profissional}
\resumeSubHeadingListStart
  \resumeSubheading{Laboratório Ferroviário (Lafer) -- parceria VALE}{Presencial}
  {Líder em Ciência de Dados}{Ago 2024 até o presente}
  \resumeItemListStart

    \resumeItem{}{Liderou uma equipe multidisciplinar de 12 pessoas, construindo sistemas de ML de manutenção preditiva de ponta a ponta para a VALE, dos requisitos das partes interessadas ao deploy escalável em produção.}

    \resumeItem{}{Reduziu a latência de detecção de falhas de 3 meses para tempo real com 91\% de precisão, usando deep learning sobre séries temporais multivariadas de sensores.}

    \resumeItem{}{Desenvolveu autoencoders Transformer não supervisionados que identificaram defeitos críticos semanas antes das manutenções programadas, incluindo modos de falha ausentes da base de manutenção.}

    %\resumeItem{}{Traduziu problemas de engenharia e manutenção em sistemas de ML em produção que informaram decisões operacionais do dia a dia.}

  \resumeItemListEnd
  \resumeSkills{Python | PyTorch | PyTorch Lightning | scikit-learn | Optuna | FastAPI | Docker | Ciência de Dados | Deep Learning | Transformers | Autoencoders | Detecção de Anomalias | Deploy}

\vspace{-1pt}

\resumeSubheading{Laboratório Ferroviário (Lafer) -- parceria VALE}{Presencial}
{Cientista de Dados Sênior}{Jan 2022 a Ago 2024}
\resumeItemListStart

  \resumeItem{}{Projetou um framework automatizado de qualidade de dados em Python que pontuava a confiabilidade de cada coleta e filtrava dados de sensores não confiáveis, substituindo a triagem manual feita viagem a viagem.}

  \resumeItem{}{Desenvolveu um autoencoder não supervisionado que detectou anomalias em sensores não identificadas por verificações baseadas em regras, com cerca de 3\% de falsos positivos, preservando a qualidade dos dados para modelagem posterior.}

  \resumeItem{}{Estabeleceu uma metodologia sistemática em Python para caracterizar o comportamento normal dos sensores, aprimorando o treinamento e a avaliação dos modelos de detecção de defeitos subsequentes.}

\resumeItemListEnd
\resumeSkills{Python | PyTorch | scikit-learn | SQL | Optuna | Ciência de Dados | Deep Learning | Autoencoders | Detecção de Anomalias | Aprendizado Não Supervisionado | Engenharia de Features | Processamento de Sinais | ETL | Qualidade de Dados}

\vspace{-1pt}

\resumeSubheading{Laboratório Ferroviário (Lafer) -- parceria VALE}{Presencial}
{Pesquisador em Ciência de Dados}{Jun 2020 a Jan 2022}
\resumeItemListStart

  \resumeItem{}{Desenvolveu um pipeline de ETL em Python que processou 20--40 mil registros de sensores por ativo/dia, integrando múltiplas fontes de dados em uma base analítica unificada que alimentou todos os modelos de ML subsequentes.}

  \resumeItem{}{Construiu modelos de deep learning que previram a condição do ativo com um R2 de aproximadamente 0,98, reduzindo a dependência de inspeções especializadas custosas.}

  \resumeItem{}{Criou a base de dados e modelagem em Python que sustentou iniciativas subsequentes de manutenção preditiva e detecção de defeitos de via na VALE.}

\resumeItemListEnd
\resumeSkills{Python | PyTorch | scikit-learn | SQL | Optuna | Ciência de Dados | Deep Learning | Engenharia de Features | Processamento de Sinais | ETL}

\vspace{-1pt}

\resumeSubheading{Laboratório de Tribologia e Dinâmica Ferroviária (LabTDF) -- parceria VALE}{Presencial}
{Estagiário em Machine Learning}{Jan 2019 a Jun 2020}
\resumeItemListStart

  \resumeItem{}{Estimou forças de contato difíceis de medir via sensores embarcados com um pipeline de ML validado por simulação, alcançando R2 de cerca de 0,97, cujos resultados ajudaram a motivar a decisão da VALE de lançar um projeto de P\&D em manutenção preditiva de via com vagões instrumentados.}

  \resumeItem{}{Otimizou perfis de desgaste roda--trilho com algoritmos genéticos, reduzindo o desgaste em 20\% e melhorando uma relação de segurança chave em cerca de 35\%, sem comprometer o desempenho à fadiga.}

  \resumeItem{}{Comprovou a necessidade de mudança no limite de manutenção via análise de desgaste/fadiga, conduzindo à adoção na operação e estendendo a vida útil do ativo em 29\%--50\% (29\% pela mudança de limite, até 50\% combinada ao perfil de roda otimizado).}

\resumeItemListEnd
\resumeSkills{Python | scikit-learn | Optuna | MATLAB | Universal Mechanism | Machine Learning | Otimização | Algoritmos Genéticos | Simulação Multicorpos}
\resumeSubHeadingListEnd

\vspace{-2pt}

%-----------KEY PROJECTS-----------------
\section{Projetos Principais}
\resumeSubHeadingListStart
  \resumeSubItem{Industrial Sensor Anomaly Detection API}{ -- \href{https://github.com/pirao/AnomalyDetection2026}{github.com/pirao/AnomalyDetection2026}}
  \resumeItemListStart
    \resumeItem{}{Serviço FastAPI em nível de produção para detecção de falhas de máquinas em tempo real, reduzindo os falsos positivos a zero em 29 cenários de benchmark (F1 0,95, recall 0,91), ante um baseline de F1 0,28.}
    \resumeItem{}{Pontuação de resíduos por grupo em sensores de vibração multivariados, com motor de alertas stateful em múltiplas janelas (confirmação e cooldown).}
    \resumeItem{}{Serve modelos versionados a partir de um registro MLflow. Containerizado com Docker Compose e CI no GitHub Actions.}
  \resumeItemListEnd
\resumeSubHeadingListEnd

\vspace{-3pt}

%-----------EDUCATION-----------------
\section{Formação Acadêmica}
  \resumeSubHeadingListStart
    \resumeSubheading
      {Universidade Estadual de Campinas (UNICAMP)}{Campinas, Brasil}
      {Doutorado em Engenharia Mecânica}{Previsão: Ago 2026}
      \resumeItemListStart
        \resumeItem{Tese: }{Uma metodologia para detecção de defeitos de via a partir de dados de veículos ferroviários instrumentados.}
      \resumeItemListEnd
    \resumeSubheading
      {Universidade Estadual de Campinas (UNICAMP)}{Campinas, Brasil}
      {Mestrado em Engenharia Mecânica}{Jan 2022}
      \resumeItemListStart
        \resumeItem{Dissertação: }{Medição de irregularidades da via a partir de dados de veículos ferroviários instrumentados via machine learning.}
      \resumeItemListEnd
    \resumeSubheading
      {Universidade Federal do Espírito Santo (UFES)}{Vitória, Brasil}
      {Bacharelado em Engenharia Mecânica}{Maio 2020}
      \resumeItemListStart
        \resumeItem{Monografia: }{Identificação indireta das forças de contato roda--trilho em veículo ferroviário heavy haul instrumentado via machine learning.}
      \resumeItemListEnd
    \resumeSubHeadingListEnd

\vspace{-2pt}

%-----------LANGUAGES-----------------
\section{Idiomas}
\vspace{1pt}
{\small Português (Nativo), Inglês (C1 - Fluente), Italiano (A2 - Básico)}

\vspace{-3pt}

%-----------PUBLICATIONS-----------------
\section{Publicações \& Conferências}
\vspace{1pt}

{\scshape Publicações Selecionadas}
\vspace{2pt}

\resumePub{2024}{Measuring vertical track irregularities from instrumented heavy haul railway vehicle data using machine learning}{Engineering Applications of Artificial Intelligence}
\resumePub{2022}{Nonintrusive Operator Inference for Chaotic Systems}{IEEE Transactions on Artificial Intelligence}
\resumePub{2021}{Indirect identification of wheel--rail contact forces of an instrumented heavy haul railway vehicle using machine learning}{Mechanical Systems and Signal Processing}
% Publicações fora de ML (engenharia mecânica) ocultadas para foco em DS/ML -- descomente para restaurar:
%\resumePub{2023}{Influence of wheel tread wear on Rolling Contact Fatigue and on the dynamics of railway vehicles}{Wear}
%\resumePub{2022}{Generation of realistic artificial track irregularities for multibody simulation using measured geometric data: from mid-chord to space-curve}{In Proceedings of the Fifth International Conference on Railway Technology}
%\resumePub{2021}{The effect of railway wheel wear on reprofiling and service life}{Wear}

\vspace{-3pt}
{\scshape Apresentações em Conferências}
\vspace{3pt}

\resumeConf{2026}{International Conference on Railway Technology (Railways 2026)}{23--26 ago. (aceito)}{Budapeste, Hungria}
\resumeConf{2025}{International Heavy Haul Association (IHHA) Conference}{17--21 nov.}{Colorado Springs, EUA}
\resumeConf{2023}{International Heavy Haul Association (IHHA) Conference}{27--31 ago.}{Rio de Janeiro, Brasil}
\resumeConf{2023}{Railway Engineering Symposium}{19--20 maio}{Campinas, Brasil}

\vspace{-5pt}

%-----------CERTIFICATIONS-----------------
% Estrutura pronta para preencher. Obs.: certificações atuais são de Jan 2024 ou anteriores.
% Descomente a seção e adicione entradas \resumePub{<ano>}{<nome da certificação>}{<emissor>}.
%\section{Certificações}
%\vspace{2pt}
%\resumePub{20XX}{Nome da Certificação}{Emissor}


\end{document}
